\section{Introduction}
\label{sec:introduction}

\subsection{The Data-Value Problem in Urban Crime Count Forecasting}
\label{subsec:intro_data_value}

Urban crime forecasting increasingly combines historical incidents with
mobility traces, calendar variables, environmental measurements, and other
dynamic urban data. These sources can enrich the representation of activity
around a location, but they also require acquisition, cleaning, spatial
alignment, maintenance, and governance. For a city agency or an analyst, the
relevant question is therefore not only whether a fitted model obtains a
lower test error after a new source is added. It is also whether that source
reduces uncertainty about future crime counts beyond information already
available from location, crime history, and the calendar.

These two questions are not equivalent. A change in held-out error depends on
the fitted model, hyperparameter search, training sample, target definition,
and validation period. It may reflect useful new information, a better match
between the added variables and one model class, or sampling variation. In
contrast, a model-independent prediction benchmark is determined by the joint
distribution of the target and the information available when a forecast is
issued. Comparing achieved error with such a benchmark can help distinguish
uncertainty intrinsic to the target--information pair from error that may
remain reducible through improved modeling.

Crime count forecasting makes this distinction practically important.
Incident counts are nonnegative integers, are often zero at fine spatial and
temporal resolutions, and are evaluated under losses such as mean squared
error (MSE). Existing crime-prediction studies have developed statistical,
machine-learning, recurrent, graph-based, multiscale, and transformer models
\citep{kounadi2020review,du2023review,mandalapu2023review,
alves2018crime,angbera2023model,huang2018deepcrime,deng2023risk,
jing2024multiscale,wang2025mragnn,butt2025start,cui2026transfer,
hu2023fusion}. Their comparative results are essential for model selection,
but they do not identify how much error is unavoidable under a specified
information set. This limitation becomes more consequential when a costly
new urban data source appears to improve one or more fitted models.

We study this problem using shared-bicycle activity as a candidate source of
additional information. Bicycle trips provide high-frequency and spatially
resolved observations of one form of public-space activity. They do not,
however, constitute a census of the ambient population, and a predictive
association with crime does not establish a causal effect. The empirical
question is deliberately narrower: does lagged bicycle flow add
prediction-relevant information for future crime counts, and can an estimated
gain be distinguished from finite-sample artifacts?

\subsection{From Integer-Valued Uncertainty to Calibrated Evidence}
\label{subsec:intro_calibrated_evidence}

Conditional entropy provides a natural measure of the uncertainty remaining
after the forecast-time information set has been observed
\citep{shannon1948,cover_thomas}. Existing entropy-based prediction bounds,
however, are mainly formulated for continuous outcomes or classification
losses. A continuous Gaussian entropy--variance argument cannot be transferred
to integer-valued crime counts by simply replacing differential entropy with
Shannon entropy. The prediction error lies on an integer lattice, and the
relevant maximum-entropy envelope under a second-moment constraint is
discrete.

To address this support mismatch, we formulate the Discrete Entropy Lower
Bound (DELB). DELB maps conditional Shannon entropy through the inverse
integer-lattice maximum-entropy envelope to a model-independent,
entropy-constrained MSE floor for causal integer-valued predictors. The bound
is exact in its discrete-envelope form, although equality with the optimal
integer-predictor risk requires the stated attainability conditions. The
framework also includes a predictor-dependent refinement and separates two
sources of nonattainment: residual dependence of the prediction error on the
available history and mismatch between the error distribution and the
moment-matched discrete Gaussian.

A population ordering is only the first layer of the applied problem. Adding
a discretized mobility variable refines the baseline state space, often
creating sparse or singleton cells. Conditional-entropy and conditional
mutual information (CMI) estimators can then have different finite-sample
errors in the baseline and augmented state tables
\citep{paninski2003,calcagnile2026benchmark}. A positive estimated CMI, or a
lower estimated DELB, may consequently arise even when the added variable has
no population information value. Moreover, a randomization reference is
conditioned on its transformation design; it need not be centered at zero
when that design changes support occupancy
\citep{ernst2004permutation,zan2022cmi,runge2018cmi,li2023nncit}.

The evaluation must therefore keep four objects separate: the population
DELB and CMI estimands, their finite-sample estimates, evidence relative to a
prespecified randomization reference, and realized out-of-sample prediction
error. We connect these layers using known-truth simulations, support
diagnostics, block resampling, temporal and spatial randomization,
circular-shift references, and matched prediction models. This design does
not turn a transformation-specific diagnostic into a causal or exact
conditional-independence test. Its purpose is to determine whether the
observed statistic is distinguishable from finite-sample behavior under each
explicit reference design.

\subsection{Research Questions and Contributions}
\label{subsec:intro_contributions}

The study addresses three linked research questions:

\begin{enumerate}
    \item \textit{How can conditional uncertainty be translated into a
    discrete-form, model-independent MSE lower bound for causal
    integer-valued prediction?}
    \item \textit{How can an estimated information gain be separated from
    sparse-state estimation bias and a transformation-specific randomization
    reference?}
    \item \textit{Does lagged bicycle mobility provide additional information
    for urban crime count prediction that remains distinguishable after
    finite-sample calibration?}
\end{enumerate}

The corresponding contributions are application oriented. First, we
introduce DELB as a model-independent benchmark for integer-valued urban count
prediction, provide its attainability diagnostics, and give an implementable
numerical inversion procedure. Second, we develop a finite-sample-calibrated
workflow that distinguishes estimated information gain, support-induced
bias, design-conditioned randomization evidence, and realized prediction
improvement. Third, we conduct a harmonized multicity evaluation using crime
and shared-bicycle data from Washington, DC, New York City, and Vancouver
during 2020--2022. The evaluation examines city-level and local estimates,
spatial and temporal scale, sample size, alternative estimators, and matched
held-out prediction.

The empirical conclusion is intentionally qualified. Adding lagged bicycle
flow lowers the estimated DELB in the primary city-level comparisons, but the
reductions are not distinguishable from the prespecified temporal, spatial,
or circular-shift references after calibration. Realized prediction changes
are model- and city-dependent. Rather than treating this result as a failed
application, the study uses it to illustrate why lower estimated
uncertainty, calibrated information gain, and improved predictive performance
must be reported as separate claims when evaluating new smart-city data.

Section~\ref{sec:related_work} reviews urban crime forecasting, mobility data
in crime analytics, and information-theoretic prediction limits.
Section~\ref{sec:theory} presents the DELB methodology. Section~\ref{sec:data}
defines the multicity data and experimental design. Section~\ref{sec:results}
first summarizes the multicity application and then reports the finite-sample,
city-level, multiscale, and prediction evidence.
Section~\ref{sec:discussion} discusses practical implications, inferential
scope, and limitations.
